\documentclass[11pt,a4paper,UTF8]{ctexart}
\usepackage{amsmath,amssymb,mathtools}
\usepackage[dvipsnames]{xcolor}
\usepackage[most]{tcolorbox}
\usepackage{geometry}
\usepackage{booktabs,longtable,array,tabularx}
\usepackage{enumitem}
\usepackage{needspace}
\usepackage[colorlinks=true,linkcolor=MidnightBlue,urlcolor=MidnightBlue]{hyperref}

\geometry{left=2.1cm,right=2.1cm,top=2.2cm,bottom=2.2cm}
\setlength{\parindent}{0em}
\setlength{\parskip}{0.35em}
\linespread{1.08}
\setlist{nosep,leftmargin=2em}
\allowdisplaybreaks
\raggedbottom

\definecolor{bluebg}{RGB}{235,245,255}
\definecolor{bluebd}{RGB}{41,128,185}
\definecolor{greenbg}{RGB}{235,255,240}
\definecolor{greenbd}{RGB}{39,174,96}
\definecolor{orangebg}{RGB}{255,246,235}
\definecolor{orangebd}{RGB}{230,126,34}
\definecolor{redbg}{RGB}{255,240,240}
\definecolor{redbd}{RGB}{192,57,43}

\newtcolorbox{identifybox}[1][]{
  enhanced,breakable,colback=bluebg,colframe=bluebd,
  title=#1,fonttitle=\bfseries,arc=2mm,boxrule=0.8pt
}
\newtcolorbox{templatebox}[1][]{
  enhanced,breakable,colback=orangebg,colframe=orangebd,
  title=#1,fonttitle=\bfseries,arc=2mm,boxrule=0.9pt
}
\newtcolorbox{answerbox}[1][]{
  enhanced,breakable,colback=greenbg,colframe=greenbd,
  title=#1,fonttitle=\bfseries,arc=2mm,boxrule=0.8pt
}
\newtcolorbox{scorebox}[1][]{
  enhanced,breakable,colback=redbg,colframe=redbd,
  title=#1,fonttitle=\bfseries,arc=2mm,boxrule=0.8pt
}

\newcommand{\rank}{\operatorname{rank}}
\newcommand{\tr}{\operatorname{tr}}
\newcommand{\diag}{\operatorname{diag}}
\newcommand{\R}{\mathbb R}

\title{\Huge 四川大学线性代数（III）期末题型考点详解\\[0.25em]
\Large 2020--2025 真题融合复习讲义}
\author{按线性代数（理工）讲义模板整理}
\date{\today}

\begin{document}
\maketitle
\thispagestyle{empty}
\tableofcontents
\newpage

\begin{scorebox}[定位]
版本：2020--2025 春季线性代数（III）期末真题融合版。\\
组织方式：不按年份堆叠，改为按题型考点分类；每道真题保留来源、题干、解答与最终答案。
\end{scorebox}

\bigskip\hrule\bigskip

\section{全题覆盖索引}

{\scriptsize
\begin{longtable}{>{\raggedright\arraybackslash}p{0.25\textwidth}>{\raggedright\arraybackslash}p{0.33\textwidth}>{\raggedright\arraybackslash}p{0.34\textwidth}}
\toprule
来源 & 考点 & 讲义位置 \\
\midrule
2020--2021 第一大题第 1 题 & 填空题 & 行列式、伴随矩阵、秩：填空题先拿分 \\
2020--2021 第一大题第 2 题 & 填空题 & 行列式、伴随矩阵、秩：填空题先拿分 \\
2020--2021 第二大题第 1 题 & 行列式 & 行列式、伴随矩阵、秩：填空题先拿分 \\
2021--2022 第一大题第 2 题 & 填空题 & 行列式、伴随矩阵、秩：填空题先拿分 \\
2021--2022 第二大题第 1 题 & 行列式与列向量变换 & 行列式、伴随矩阵、秩：填空题先拿分 \\
2022--2023 第一大题第 1 题 & 填空题 & 行列式、伴随矩阵、秩：填空题先拿分 \\
2022--2023 第一大题第 2 题 & 填空题 & 行列式、伴随矩阵、秩：填空题先拿分 \\
2022--2023 第二大题第 1 题 & 余子式之和 & 行列式、伴随矩阵、秩：填空题先拿分 \\
2023--2024 第二大题第 1 题 & 余子式组合 & 行列式、伴随矩阵、秩：填空题先拿分 \\
2024--2025 第一大题第 5 题 & 填空题 & 行列式、伴随矩阵、秩：填空题先拿分 \\
2024--2025 第二大题第 1 题 & 行列式与代数余子式 & 行列式、伴随矩阵、秩：填空题先拿分 \\
2020--2021 第一大题第 3 题 & 填空题 & 矩阵方程与矩阵幂 \\
2020--2021 第二大题第 3 题 & 矩阵方程 & 矩阵方程与矩阵幂 \\
2021--2022 第一大题第 1 题 & 填空题 & 矩阵方程与矩阵幂 \\
2021--2022 第二大题第 2 题 & 矩阵方程 & 矩阵方程与矩阵幂 \\
2022--2023 第一大题第 3 题 & 填空题 & 矩阵方程与矩阵幂 \\
2022--2023 第一大题第 4 题 & 填空题 & 矩阵方程与矩阵幂 \\
2022--2023 第二大题第 2 题 & 矩阵方程 & 矩阵方程与矩阵幂 \\
2023--2024 第一大题第 2 题 & 填空题 & 矩阵方程与矩阵幂 \\
2023--2024 第一大题第 5 题 & 填空题 & 矩阵方程与矩阵幂 \\
2023--2024 第二大题第 2 题 & 矩阵方程 & 矩阵方程与矩阵幂 \\
2024--2025 第一大题第 1 题 & 填空题 & 矩阵方程与矩阵幂 \\
2024--2025 第一大题第 4 题 & 填空题 & 矩阵方程与矩阵幂 \\
2024--2025 第二大题第 2 题 & 矩阵方程 & 矩阵方程与矩阵幂 \\
2020--2021 第三大题第 1 题 & 向量组 & 向量组、秩、极大线性无关组 \\
2021--2022 第一大题第 3 题 & 填空题 & 向量组、秩、极大线性无关组 \\
2021--2022 第一大题第 5 题 & 填空题 & 向量组、秩、极大线性无关组 \\
2021--2022 第二大题第 3 题 & 向量组秩与极大无关组 & 向量组、秩、极大线性无关组 \\
2022--2023 第二大题第 3 题 & 向量组秩与线性表出 & 向量组、秩、极大线性无关组 \\
2023--2024 第一大题第 1 题 & 填空题 & 向量组、秩、极大线性无关组 \\
2023--2024 第二大题第 3 题 & 全部极大无关组 & 向量组、秩、极大线性无关组 \\
2024--2025 第一大题第 2 题 & 填空题 & 向量组、秩、极大线性无关组 \\
2024--2025 第二大题第 3 题 & 向量组秩与参数 & 向量组、秩、极大线性无关组 \\
2020--2021 第一大题第 5 题 & 填空题 & 线性方程组：有解、无穷多解、通解 \\
2020--2021 第二大题第 2 题 & 齐次方程组 & 线性方程组：有解、无穷多解、通解 \\
2021--2022 第三大题第 1 题 & 含参数线性方程组 & 线性方程组：有解、无穷多解、通解 \\
2022--2023 第一大题第 5 题 & 填空题 & 线性方程组：有解、无穷多解、通解 \\
2022--2023 第三大题第 1 题 & 含参数线性方程组 & 线性方程组：有解、无穷多解、通解 \\
2023--2024 第一大题第 4 题 & 填空题 & 线性方程组：有解、无穷多解、通解 \\
2023--2024 第三大题第 1 题 & 参数齐次线性方程组 & 线性方程组：有解、无穷多解、通解 \\
2024--2025 第三大题第 1 题 & 非齐次线性方程组 & 线性方程组：有解、无穷多解、通解 \\
2020--2021 第一大题第 4 题 & 填空题 & 特征值、特征向量、相似对角化、实对称矩阵 \\
2020--2021 第三大题第 2 题 & 实对称矩阵 & 特征值、特征向量、相似对角化、实对称矩阵 \\
2021--2022 第一大题第 4 题 & 填空题 & 特征值、特征向量、相似对角化、实对称矩阵 \\
2021--2022 第三大题第 2 题 & 实对称矩阵 & 特征值、特征向量、相似对角化、实对称矩阵 \\
2022--2023 第三大题第 2 题 & 特征值、对角化与矩阵幂 & 特征值、特征向量、相似对角化、实对称矩阵 \\
2023--2024 第一大题第 3 题 & 填空题 & 特征值、特征向量、相似对角化、实对称矩阵 \\
2023--2024 第一大题第 6 题 & 填空题 & 特征值、特征向量、相似对角化、实对称矩阵 \\
2023--2024 第三大题第 2 题 & 特征向量参数与对角化 & 特征值、特征向量、相似对角化、实对称矩阵 \\
2024--2025 第一大题第 6 题 & 填空题 & 特征值、特征向量、相似对角化、实对称矩阵 \\
2024--2025 第三大题第 3 题 & 特征多项式与向量方程 & 特征值、特征向量、相似对角化、实对称矩阵 \\
2020--2021 第一大题第 6 题 & 填空题 & 二次型：正定、标准形、正交变换 \\
2020--2021 第三大题第 3 题 & 二次型正交化 & 二次型：正定、标准形、正交变换 \\
2021--2022 第一大题第 6 题 & 填空题 & 二次型：正定、标准形、正交变换 \\
2021--2022 第三大题第 3 题 & 二次型正交化 & 二次型：正定、标准形、正交变换 \\
2022--2023 第一大题第 6 题 & 填空题 & 二次型：正定、标准形、正交变换 \\
2022--2023 第三大题第 3 题 & 二次型正交变换 & 二次型：正定、标准形、正交变换 \\
2023--2024 第三大题第 3 题 & 二次型正交化 & 二次型：正定、标准形、正交变换 \\
2024--2025 第一大题第 3 题 & 填空题 & 二次型：正定、标准形、正交变换 \\
2024--2025 第三大题第 2 题 & 二次型正交化 & 二次型：正定、标准形、正交变换 \\
2020--2021 第四大题第 1 题 & 证明题 & 证明题：固定套路 \\
2020--2021 第四大题第 2 题 & 证明题 & 证明题：固定套路 \\
2021--2022 第四大题第 1 题 & 证明题 & 证明题：固定套路 \\
2021--2022 第四大题第 2 题 & 证明题 & 证明题：固定套路 \\
2022--2023 第四大题第 1 题 & 证明题 & 证明题：固定套路 \\
2022--2023 第四大题第 2 题 & 证明题 & 证明题：固定套路 \\
2023--2024 第四大题第 1 题 & 证明题 & 证明题：固定套路 \\
2023--2024 第四大题第 2 题 & 证明题 & 证明题：固定套路 \\
2024--2025 第四大题第 1 题 & 证明题 & 证明题：固定套路 \\
2024--2025 第四大题第 2 题 & 证明题 & 证明题：固定套路 \\
\bottomrule
\end{longtable}
}

\bigskip\hrule\bigskip

\section{高频公式与考试模板}

\begin{identifybox}[核心公式]
\[
AA^*=A^*A=|A|E,\qquad |A^*|=|A|^{n-1},\qquad A^*=|A|A^{-1}\ (A\text{ 可逆}).
\]
\[
\rank(AB)\le \min\{\rank(A),\rank(B)\},\qquad
P,Q\text{ 可逆}\Rightarrow \rank(PAQ)=\rank(A).
\]
\[
P^{-1}AP=\Lambda,\qquad Q^TAQ=\Lambda\quad(A=A^T).
\]
\end{identifybox}

\begin{templatebox}[大题给分模板]
向量组：构造 $A=(\alpha_1,\ldots,\alpha_s)$，行化简，写主元列、秩、极大无关组，主元列对应原向量。\\
方程组：写 $(A|b)$，求 $r(A),r(A|b)$，判断解，通解写 $X=X_0+k_1\xi_1+\cdots$。\\
二次型：写对称矩阵，求特征值和单位正交特征向量，构造 $Q$，写 $Q^TAQ=\Lambda$ 与标准形。
\end{templatebox}

\section{行列式、伴随矩阵、秩：填空题先拿分}

\begin{identifybox}[识别]
看到 $|A|,A^*,A^{-1},r(A)$，先用伴随矩阵公式、初等变换不改变秩、可逆矩阵左右乘不改变秩。
\end{identifybox}
\begin{templatebox}[考场模板]
$AA^*=A^*A=|A|E$，$|A^*|=|A|^{n-1}$；行列式题优先做行列变换，秩题优先化阶梯形。
\end{templatebox}

\subsection{真题融合}

\begin{answerbox}[2020--2021 第一大题第 1 题：填空题]
\textbf{【来源】} 2020--2021 第一大题第 1 题

\textbf{【原题】}
设 $B$ 为 $3$ 阶方阵，$I_3$ 为 $3$ 阶单位矩阵。假定 $B$ 的特征值分别为 $1,2,3$，求 $\left|B^2-B+I_3\right|$。

\textbf{【速解】}
$B^2-B+I$ 的特征值为 $\lambda_i^2-\lambda_i+1$，故结果为 $1\cdot3\cdot7=21$。

\textbf{【答案】}
\[21\]
\end{answerbox}

\begin{answerbox}[2020--2021 第一大题第 2 题：填空题]
\textbf{【来源】} 2020--2021 第一大题第 2 题

\textbf{【原题】}
已知矩阵
\[
A=\begin{pmatrix}1&1&2\\0&2&3\\2&k&7\end{pmatrix}
\]
的秩为 $2$，求 $k$。

\textbf{【速解】}
$|A|=0$ 得 $k=4$。

\textbf{【答案】}
\[k=4\]
\end{answerbox}

\begin{answerbox}[2020--2021 第二大题第 1 题：行列式]
\textbf{【题目分析】} 计算
\[
D=\begin{vmatrix}
b&2&2&2\\
2&b&2&2\\
2&2&b&2\\
2&2&2&b
\end{vmatrix}.
\]

\textbf{【详细解答】} 该矩阵为 $(b-2)I+2J$，其中 $J$ 为全 $1$ 矩阵。$J$ 的特征值为 $4,0,0,0$，故原矩阵特征值为 $b+6,b-2,b-2,b-2$。

\textbf{【最终答案】}
\[
D=(b+6)(b-2)^3.
\]

\textbf{【方法总结】} 对角线相同、非对角线相同的行列式，直接用 $aI+cJ$ 的特征值。
\end{answerbox}

\begin{answerbox}[2021--2022 第一大题第 2 题：填空题]
\textbf{【来源】} 2021--2022 第一大题第 2 题

\textbf{【原题】}
$A,B$ 为相似 $4$ 阶方阵，$A$ 的特征值为 $1,2,3,4$，求 $|B^2+B|$。

\textbf{【速解】}
$B$ 与 $A$ 相似，$B^2+B$ 的特征值为 $\lambda_i^2+\lambda_i$。

\textbf{【答案】}
\[|B^2+B|=2880\]
\end{answerbox}

\begin{answerbox}[2021--2022 第二大题第 1 题：行列式与列向量变换]
\textbf{【题目分析】} 已知
\[
\beta_1=\alpha_1+\alpha_2+\alpha_3,\quad
\beta_2=\alpha_1+2\alpha_2+4\alpha_3,\quad
\beta_3=\alpha_1+3\alpha_2+9\alpha_3,
\]
\[
A=(\alpha_1,\alpha_2,\alpha_3),\quad B=(\beta_1,\beta_2,\beta_3),\quad |A|=1.
\]
求 $|B|$。

\textbf{【详细解答】}
\[
B=A\begin{pmatrix}1&1&1\\1&2&3\\1&4&9\end{pmatrix},
\quad |B|=|A|\begin{vmatrix}1&1&1\\1&2&3\\1&4&9\end{vmatrix}=2.
\]

\textbf{【最终答案】}
\[
|B|=2.
\]

\textbf{【方法总结】} 新列向量由旧列向量线性表示时，写成 $B=AC$，再用 $|B|=|A||C|$。
\end{answerbox}

\begin{answerbox}[2022--2023 第一大题第 1 题：填空题]
\textbf{【来源】} 2022--2023 第一大题第 1 题

\textbf{【原题】}
设
\[
A=\begin{pmatrix}1&-2&2\\4&3&t\\3&1&-1\end{pmatrix},
\]
$B$ 为 $3$ 阶非零方阵，且满足 $AB=0$，求 $t$。

\textbf{【速解】}
$B\ne O$ 且 $AB=0$，故 $|A|=0$。

\textbf{【答案】}
\[t=-3\]
\end{answerbox}

\begin{answerbox}[2022--2023 第一大题第 2 题：填空题]
\textbf{【来源】} 2022--2023 第一大题第 2 题

\textbf{【原题】}
$A$ 为 $5$ 阶方阵且 $\rank(A)=3$，求 $\rank(A^*)$。

\textbf{【速解】}
$n=5,\ r(A)=3<n-1$，故 $r(A^*)=0$。

\textbf{【答案】}
\[\rank(A^*)=0\]
\end{answerbox}

\begin{answerbox}[2022--2023 第二大题第 1 题：余子式之和]
\textbf{【题目分析】} 设行列式
\[
D=\begin{vmatrix}3&0&4&0\\2&2&2&2\\0&-7&0&0\\1&2&3&4\end{vmatrix},
\]
求 $D$ 的第 $4$ 行各元素的余子式之和。

\textbf{【详细解答】} 记第 $4$ 行第 $j$ 个元素的余子式为 $M_{4j}$。由余子式与代数余子式的符号关系，
\[
M_{41}+M_{42}+M_{43}+M_{44}
=-A_{41}+A_{42}-A_{43}+A_{44}.
\]
按第 $4$ 行代数余子式组合计算，得
\[
M_{41}+M_{42}+M_{43}+M_{44}=-28.
\]

\textbf{【最终答案】}
\[
M_{41}+M_{42}+M_{43}+M_{44}=-28.
\]

\textbf{【方法总结】} 余子式求和要先区分余子式 $M_{ij}$ 与代数余子式 $A_{ij}$ 的符号。
\end{answerbox}

\begin{answerbox}[2023--2024 第二大题第 1 题：余子式组合]
\textbf{【题目分析】} 设
\[
A=\begin{pmatrix}1&2&3&4\\0&3&4&6\\0&4&1&2\\0&2&2&2\end{pmatrix},
\]
记第 $i$ 行第 $j$ 列元素的余子式为 $M_{ij}$，求
\[
M_{11}-2M_{21}+M_{31}-2M_{41}.
\]

\textbf{【详细解答】} 由余子式与代数余子式关系，
\[
M_{11}-2M_{21}+M_{31}-2M_{41}
=A_{11}+2A_{21}+A_{31}+2A_{41}.
\]
该式等于把 $A$ 第一列替换为 $(1,2,1,2)^T$ 后所得行列式：
\[
\begin{vmatrix}
1&2&3&4\\
2&3&4&6\\
1&4&1&2\\
2&2&2&2
\end{vmatrix}
=-12.
\]

\textbf{【最终答案】}
\[
M_{11}-2M_{21}+M_{31}-2M_{41}=-12.
\]

\textbf{【方法总结】} 代数余子式的线性组合可转化为替换对应列后的行列式。
\end{answerbox}

\begin{answerbox}[2024--2025 第一大题第 5 题：填空题]
\textbf{【来源】} 2024--2025 第一大题第 5 题

\textbf{【原题】}
设 $3$ 阶方阵 $A$ 的特征值为 $1,-1,2$，$I$ 为 $3$ 阶单位矩阵，求 $|A^2+2I|$。

\textbf{【速解】}
$A^2+2I$ 的特征值为 $\lambda_i^2+2$。

\textbf{【答案】}
\[|A^2+2I|=54\]
\end{answerbox}

\begin{answerbox}[2024--2025 第二大题第 1 题：行列式与代数余子式]
\textbf{【题目分析】} 已知
\[
A=\begin{pmatrix}1&2&1&3\\2&1&0&1\\0&3&2&5\\1&0&1&2\end{pmatrix}.
\]
$(1)$ 求 $|A|$；$(2)$ 记矩阵 $A$ 中第 $i$ 行第 $j$ 列元素 $a_{ij}$ 的代数余子式为 $A_{ij}$，求
\[
A_{11}+2A_{21}+3A_{31}+A_{41}.
\]

\textbf{【详细解答】} 对 $A$ 作初等行变换可得
\[
|A|=0.
\]
代数余子式线性组合可化为替换第一列后的行列式：
\[
A_{11}+2A_{21}+3A_{31}+A_{41}=-3.
\]

\textbf{【最终答案】}
\[
|A|=0,\qquad A_{11}+2A_{21}+3A_{31}+A_{41}=-3.
\]

\textbf{【方法总结】} 代数余子式组合常转为替换列后的行列式计算。
\end{answerbox}


\section{矩阵方程与矩阵幂}

\begin{identifybox}[识别]
矩阵方程先整理未知矩阵所在位置，再决定左乘或右乘逆矩阵；矩阵幂优先用特征值、最小多项式或递推关系。
\end{identifybox}
\begin{templatebox}[考场模板]
保持乘法顺序：$AX=B\Rightarrow X=A^{-1}B$，$XA=B\Rightarrow X=BA^{-1}$。伴随矩阵统一使用 $AA^*=|A|E$。
\end{templatebox}

\subsection{真题融合}

\begin{answerbox}[2020--2021 第一大题第 3 题：填空题]
\textbf{【来源】} 2020--2021 第一大题第 3 题

\textbf{【原题】}
已知
\[
A=\begin{pmatrix}0&2&1\\0&5&3\\2&0&0\end{pmatrix},
\]
求 $A^{-1}$。

\textbf{【速解】}
行化简 $(A|I)$ 得逆矩阵。

\textbf{【答案】}
\[A^{-1}=\begin{pmatrix}0&0&\frac12\\3&-1&0\\-5&2&0\end{pmatrix}\]
\end{answerbox}

\begin{answerbox}[2020--2021 第二大题第 3 题：矩阵方程]
\textbf{【题目分析】} 已知 $AX=B$，其中
\[
A=\begin{pmatrix}1&0&-2\\-3&-1&1\\-2&-1&0\end{pmatrix},\quad
B=\begin{pmatrix}1&-1\\2&0\\5&-3\end{pmatrix}.
\]
求 $X$。

\textbf{【详细解答】} 因 $A$ 可逆，
\[
X=A^{-1}B
=\begin{pmatrix}5&-5\\-15&13\\2&-2\end{pmatrix}.
\]
等价于对增广矩阵 $(A|B)$ 作初等行变换。

\textbf{【最终答案】}
\[
X=\begin{pmatrix}5&-5\\-15&13\\2&-2\end{pmatrix}.
\]
\textbf{【方法总结】} 矩阵方程 $AX=B$ 只能左乘 $A^{-1}$。
\end{answerbox}

\begin{answerbox}[2021--2022 第一大题第 1 题：填空题]
\textbf{【来源】} 2021--2022 第一大题第 1 题

\textbf{【原题】}
$A$ 为 $3$ 阶方阵，$A^2+A-2I_3=0$，求 $(A-I_3)^{-1}$。

\textbf{【速解】}
由 $A^2+A-2I_3=0$ 得 $(A-I_3)(A+2I_3)=I_3$。

\textbf{【答案】}
\[(A-I_3)^{-1}=A+2I_3\]
\end{answerbox}

\begin{answerbox}[2021--2022 第二大题第 2 题：矩阵方程]
\textbf{【题目分析】} 已知
\[
A=\begin{pmatrix}1&-1&0\\0&1&-2\\0&0&1\end{pmatrix},\quad
B=\begin{pmatrix}1&-1\\0&2\\-3&1\end{pmatrix},
\]
求满足 $AX=B$ 的矩阵 $X$。

\textbf{【详细解答】} 由于 $|A|=1\ne0$，故 $A$ 可逆。由 $AX=B$ 左乘 $A^{-1}$ 得
\[
X=A^{-1}B,\quad
A^{-1}=\begin{pmatrix}1&1&2\\0&1&2\\0&0&1\end{pmatrix}.
\]
于是
\[
X=\begin{pmatrix}1&1&2\\0&1&2\\0&0&1\end{pmatrix}
\begin{pmatrix}1&-1\\0&2\\-3&1\end{pmatrix}
=\begin{pmatrix}-5&3\\-6&4\\-3&1\end{pmatrix}.
\]

\textbf{【最终答案】}
\[
X=\begin{pmatrix}-5&3\\-6&4\\-3&1\end{pmatrix}.
\]

\textbf{【方法总结】} $AX=B$ 只能左乘 $A^{-1}$，不能交换矩阵乘法顺序。
\end{answerbox}

\begin{answerbox}[2022--2023 第一大题第 3 题：填空题]
\textbf{【来源】} 2022--2023 第一大题第 3 题

\textbf{【原题】}
$3$ 阶方阵 $A$ 的特征值为 $1,2,2$，求 $|4A^{-1}+I_3|$。

\textbf{【速解】}
$4A^{-1}+I_3$ 的特征值为 $4/\lambda_i+1$。

\textbf{【答案】}
\[|4A^{-1}+I|=45\]
\end{answerbox}

\begin{answerbox}[2022--2023 第一大题第 4 题：填空题]
\textbf{【来源】} 2022--2023 第一大题第 4 题

\textbf{【原题】}
$A=\begin{pmatrix}2&1&0\\3&2&0\\0&0&1\end{pmatrix}$，求 $A^{-1}$。

\textbf{【速解】}
直接求逆。

\textbf{【答案】}
\[A^{-1}=\begin{pmatrix}2&-1&0\\-3&2&0\\0&0&1\end{pmatrix}\]
\end{answerbox}

\begin{answerbox}[2022--2023 第二大题第 2 题：矩阵方程]
\textbf{【题目分析】} 已知
\[
A=\begin{pmatrix}4&2&3\\1&1&0\\-1&2&3\end{pmatrix},
\]
且有 $AB-2B=A$，求矩阵 $B$。

\textbf{【详细解答】} 先整理矩阵方程：
\[
AB-2B=(A-2I_3)B=A.
\]
因 $|A-2I_3|\ne0$，故可左乘 $(A-2I_3)^{-1}$：
\[
B=(A-2I_3)^{-1}A.
\]
对增广矩阵 $(A-2I_3\mid A)$ 行化简：
\[
\left(
\begin{array}{ccc|ccc}
2&2&3&4&2&3\\
1&-1&0&1&1&0\\
-1&2&1&-1&2&3
\end{array}
\right)
\to
\left(
\begin{array}{ccc|ccc}
1&0&0&3&-8&-6\\
0&1&0&2&-9&-6\\
0&0&1&-2&12&9
\end{array}
\right).
\]

\textbf{【最终答案】}
\[
B=
\begin{pmatrix}
3&-8&-6\\
2&-9&-6\\
-2&12&9
\end{pmatrix}.
\]

\textbf{【方法总结】} 矩阵方程先提公共右因子 $B$，再左乘逆矩阵，乘法顺序不能交换。
\end{answerbox}

\begin{answerbox}[2023--2024 第一大题第 2 题：填空题]
\textbf{【来源】} 2023--2024 第一大题第 2 题

\textbf{【原题】}
已知
\[
A=\begin{pmatrix}
0&1&0&0\\
-1&0&0&0\\
0&0&-1&-1\\
0&0&0&-1
\end{pmatrix},
\]
求 $A^{-1}$。

\textbf{【速解】}
分块求逆。

\textbf{【答案】}
\[A^{-1}=\begin{pmatrix}0&-1&0&0\\1&0&0&0\\0&0&-1&1\\0&0&0&-1\end{pmatrix}\]
\end{answerbox}

\begin{answerbox}[2023--2024 第一大题第 5 题：填空题]
\textbf{【来源】} 2023--2024 第一大题第 5 题

\textbf{【原题】}
设 $3$ 阶矩阵 $A$ 的行列式 $|A|=2$，又已知 $\lambda$ 是 $A$ 的一个特征值，求 $A^*$ 的其中一个对应特征值。

\textbf{【速解】}
若 $Ax=\lambda x$，则 $A^*x=\frac{|A|}{\lambda}x$。

\textbf{【答案】}
\[\frac{2}{\lambda}\]
\end{answerbox}

\begin{answerbox}[2023--2024 第二大题第 2 题：矩阵方程]
\textbf{【题目分析】} 已知 $A,B$ 为 $3$ 阶矩阵，满足
\[
2A^{-1}B=B-4I,
\quad
B=\begin{pmatrix}1&-2&0\\1&2&0\\0&0&2\end{pmatrix}.
\]
求矩阵 $A$。

\textbf{【详细解答】} 由
\[
2A^{-1}B=B-4I
\]
左乘 $A$ 得
\[
2B=A(B-4I).
\]
由于
\[
B-4I=\begin{pmatrix}-3&-2&0\\1&-2&0\\0&0&-2\end{pmatrix}
\]
可逆，故右乘 $(B-4I)^{-1}$：
\[
A=2B(B-4I)^{-1}
=\begin{pmatrix}0&2&0\\-1&-1&0\\0&0&-2\end{pmatrix}.
\]

\textbf{【最终答案】}
\[
A=\begin{pmatrix}0&2&0\\-1&-1&0\\0&0&-2\end{pmatrix}.
\]

\textbf{【方法总结】} 方程含 $A^{-1}$ 时先左乘 $A$，再保持右乘逆矩阵的位置。
\end{answerbox}

\begin{answerbox}[2024--2025 第一大题第 1 题：填空题]
\textbf{【来源】} 2024--2025 第一大题第 1 题

\textbf{【原题】}
设矩阵
\[
A=\begin{pmatrix}1&0&2\\0&1&0\\0&0&1\end{pmatrix},
\]
求 $A^{-1}$。

\textbf{【速解】}
上三角初等矩阵直接求逆。

\textbf{【答案】}
\[A^{-1}=\begin{pmatrix}1&0&-2\\0&1&0\\0&0&1\end{pmatrix}\]
\end{answerbox}

\begin{answerbox}[2024--2025 第一大题第 4 题：填空题]
\textbf{【来源】} 2024--2025 第一大题第 4 题

\textbf{【原题】}
设 $4$ 阶矩阵 $A$ 的秩为 $3$，求其伴随矩阵 $A^*$ 的秩。

\textbf{【速解】}
$n$ 阶矩阵秩为 $n-1$ 时，伴随矩阵秩为 $1$。

\textbf{【答案】}
\[\rank(A^*)=1\]
\end{answerbox}

\begin{answerbox}[2024--2025 第二大题第 2 题：矩阵方程]
\textbf{【题目分析】} 已知
\[
A=\begin{pmatrix}1&2&1\\0&1&2\\2&1&-3\end{pmatrix},
\quad
B=\begin{pmatrix}1&0\\2&1\\1&3\end{pmatrix},
\]
求满足 $XA=B^T$ 的矩阵 $X$。

\textbf{【详细解答】} 方程 $XA=B^T$ 中未知矩阵在左侧，右乘 $A^{-1}$ 得
\[
X=B^TA^{-1}.
\]
也可转置为
\[
A^TX^T=B,
\]
解得
\[
X=\begin{pmatrix}1&0&0\\-2&4&1\end{pmatrix}.
\]

\textbf{【最终答案】}
\[
X=\begin{pmatrix}1&0&0\\-2&4&1\end{pmatrix}.
\]

\textbf{【方法总结】} $XA=B^T$ 应右乘 $A^{-1}$；转置法可化为常见的 $A^TX^T=B$。
\end{answerbox}


\section{向量组、秩、极大线性无关组}

\begin{identifybox}[识别]
向量组题统一把向量作为列向量构造矩阵，行化简后看主元列、秩和线性表出。
\end{identifybox}
\begin{templatebox}[考场模板]
$A=(\alpha_1,\alpha_2,\cdots)$；主元列对应原矩阵中的列向量，不是化简后矩阵中的列。
\end{templatebox}

\subsection{真题融合}

\begin{answerbox}[2020--2021 第三大题第 1 题：向量组]
\textbf{【题目分析】} 设
\[
\alpha_1=(1,2,-1,-1)^T,\ \alpha_2=(-3,-1,2,0)^T,
\]
\[
\alpha_3=(2,-1,-1,1)^T,\ \alpha_4=(-1,-2,1,1)^T.
\]
求秩、一个极大线性无关组，并将其余向量线性表出。

\textbf{【详细解答】} 构造
\[
A=(\alpha_1,\alpha_2,\alpha_3,\alpha_4).
\]
行化简得
\[
A\to
\begin{pmatrix}
1&0&-1&-1\\
0&1&-1&0\\
0&0&0&0\\
0&0&0&0
\end{pmatrix}.
\]
主元列为第 $1,2$ 列，对应原向量中的 $\alpha_1,\alpha_2$，故
\[
\rank(A)=2,\qquad \alpha_1,\alpha_2\text{ 为一个极大线性无关组}.
\]
由非主元列关系得
\[
\alpha_3=-\alpha_1-\alpha_2,\qquad \alpha_4=-\alpha_1.
\]

\textbf{【最终答案】}
\[
\rank(A)=2,\quad \{\alpha_1,\alpha_2\}\text{ 为一个极大线性无关组},
\]
\[
\alpha_3=-\alpha_1-\alpha_2,\qquad \alpha_4=-\alpha_1.
\]

\textbf{【方法总结】} 主元列必须回到原列向量，不取阶梯形列。
\end{answerbox}

\begin{answerbox}[2021--2022 第一大题第 3 题：填空题]
\textbf{【来源】} 2021--2022 第一大题第 3 题

\textbf{【原题】}
求二次型 $f=(x_1-x_2)^2-(x_1-x_3)^2+(x_2-x_3)^2$ 的秩。

\textbf{【速解】}
写出二次型矩阵并求秩。

\textbf{【答案】}
\[\rank(f)=2\]
\end{answerbox}

\begin{answerbox}[2021--2022 第一大题第 5 题：填空题]
\textbf{【来源】} 2021--2022 第一大题第 5 题

\textbf{【原题】}
$\alpha_1,\alpha_2,\alpha_3,\alpha_4,\beta$ 为给定 $4$ 维列向量组，且 $\beta$ 不能由 $\alpha_1,\alpha_2,\alpha_3,\alpha_4$ 线性表示，判断 $\alpha_1,\alpha_2,\alpha_3,\alpha_4,\beta$ 线性相关或无关。

\textbf{【速解】}
$5$ 个 $4$ 维向量必线性相关。

\textbf{【答案】}
\[\alpha_1,\alpha_2,\alpha_3,\alpha_4,\beta\text{ 线性相关}\]
\end{answerbox}

\begin{answerbox}[2021--2022 第二大题第 3 题：向量组秩与极大无关组]
\textbf{【题目分析】} 给定
\[
\alpha_1=(1,-1,2,4)^T,\quad
\alpha_2=(0,3,1,2)^T,
\]
\[
\alpha_3=(3,0,7,14)^T,\quad
\alpha_4=(1,-2,2,0)^T.
\]
求向量组的秩、一个极大线性无关组，并表示其余向量。

\textbf{【详细解答】} 构造列向量矩阵
\[
C=(\alpha_1,\alpha_2,\alpha_3,\alpha_4)
=\begin{pmatrix}
1&0&3&1\\
-1&3&0&-2\\
2&1&7&2\\
4&2&14&0
\end{pmatrix}.
\]
初等行变换得
\[
C\to
\begin{pmatrix}
1&0&3&0\\
0&1&1&0\\
0&0&0&1\\
0&0&0&0
\end{pmatrix}.
\]
主元列为第 $1,2,4$ 列，对应原矩阵中的 $\alpha_1,\alpha_2,\alpha_4$。由第 $3$ 列关系得
\[
\alpha_3=3\alpha_1+\alpha_2.
\]

\textbf{【最终答案】}
\[
\rank(\alpha_1,\alpha_2,\alpha_3,\alpha_4)=3,
\quad \alpha_1,\alpha_2,\alpha_4\text{ 为一个极大线性无关组},
\quad \alpha_3=3\alpha_1+\alpha_2.
\]

\textbf{【方法总结】} 主元列必须对应原列向量；被表出向量由行化简后的非主元列关系确定。
\end{answerbox}

\begin{answerbox}[2022--2023 第二大题第 3 题：向量组秩与线性表出]
\textbf{【题目分析】} 设
\[
\alpha_1=(1,-2,0,3)^T,\quad
\alpha_2=(2,-5,-3,6)^T,
\]
\[
\alpha_3=(0,1,3,0)^T,\quad
\alpha_4=(2,-1,4,-7)^T.
\]
求该向量组的秩和一个极大线性无关组，并将其余向量用该极大线性无关组线性表示。

\textbf{【详细解答】} 构造列向量矩阵
\[
C=(\alpha_1,\alpha_2,\alpha_3,\alpha_4)
=\begin{pmatrix}
1&2&0&2\\
-2&-5&1&-1\\
0&-3&3&4\\
3&6&0&-7
\end{pmatrix}.
\]
初等行变换得
\[
C\to
\begin{pmatrix}
1&0&2&0\\
0&1&-1&0\\
0&0&0&1\\
0&0&0&0
\end{pmatrix}.
\]
主元列为第 $1,2,4$ 列，对应原矩阵中的 $\alpha_1,\alpha_2,\alpha_4$。由第 $3$ 列关系得
\[
\alpha_3=2\alpha_1-\alpha_2.
\]

\textbf{【最终答案】}
\[
\rank(\alpha_1,\alpha_2,\alpha_3,\alpha_4)=3,
\quad \alpha_1,\alpha_2,\alpha_4\text{ 为一个极大线性无关组},
\]
\[
\alpha_3=2\alpha_1-\alpha_2.
\]

\textbf{【方法总结】} 向量组题必须写列矩阵和主元列；主元列回到原向量组中读取。
\end{answerbox}

\begin{answerbox}[2023--2024 第一大题第 1 题：填空题]
\textbf{【来源】} 2023--2024 第一大题第 1 题

\textbf{【原题】}
已知
\[
A=\begin{pmatrix}-2&0&-2\\0&-1&2\\-3&-2&1\end{pmatrix},
\]
求 $r(A)$。

\textbf{【速解】}
行化简求秩。

\textbf{【答案】}
\[r(A)=2\]
\end{answerbox}

\begin{answerbox}[2023--2024 第二大题第 3 题：全部极大无关组]
\textbf{【题目分析】} 求向量组
\[
\alpha_1=(2,3,3,2)^T,\quad
\alpha_2=(2,1,4,0)^T,
\]
\[
\alpha_3=(1,2,2,0)^T,\quad
\alpha_4=(3,1,4,4)^T
\]
的秩和全部极大线性无关组。

\textbf{【详细解答】} 构造列向量矩阵
\[
C=(\alpha_1,\alpha_2,\alpha_3,\alpha_4)
=\begin{pmatrix}
2&2&1&3\\
3&1&2&1\\
3&4&2&4\\
2&0&0&4
\end{pmatrix}.
\]
行化简得
\[
C\to
\begin{pmatrix}
-1&1&-1&2\\
0&4&-1&7\\
0&0&1&-3\\
0&0&0&0
\end{pmatrix}.
\]
故 $\rank(C)=3$。逐一验证任取三个向量所得四个三元组的秩均为 $3$。

\textbf{【最终答案】}
\[
\rank(C)=3,
\]
全部极大线性无关组为
\[
\{\alpha_1,\alpha_2,\alpha_3\},\quad
\{\alpha_1,\alpha_2,\alpha_4\},\quad
\{\alpha_1,\alpha_3,\alpha_4\},\quad
\{\alpha_2,\alpha_3,\alpha_4\}.
\]

\textbf{【方法总结】} 题目问“全部”时，不能只给主元列组；需列出所有秩为 $r$ 的子组。
\end{answerbox}

\begin{answerbox}[2024--2025 第一大题第 2 题：填空题]
\textbf{【来源】} 2024--2025 第一大题第 2 题

\textbf{【原题】}
求向量组
\[
\alpha_1=(1,2,1)^T,\quad \alpha_2=(2,1,3)^T,\quad \alpha_3=(1,-1,2)^T
\]
的秩。

\textbf{【速解】}
构造列矩阵求秩。

\textbf{【答案】}
\[r(\alpha_1,\alpha_2,\alpha_3)=2\]
\end{answerbox}

\begin{answerbox}[2024--2025 第二大题第 3 题：向量组秩与参数]
\textbf{【题目分析】} 已知向量组
\[
\alpha_1=(1,1,0,2)^T,\quad
\alpha_2=(2,2,1,5)^T,
\]
\[
\alpha_3=(-1,-1,2,1)^T,\quad
\alpha_4=(0,a,1,1)^T
\]
的秩为 $3$，其中 $a$ 为未知参数。求 $a$ 的值和该向量组的一个极大线性无关组，并将其余向量用所求极大线性无关组线性表示。

\textbf{【详细解答】} 构造列向量矩阵
\[
C=(\alpha_1,\alpha_2,\alpha_3,\alpha_4).
\]
由行化简和秩为 $3$ 的条件得
\[
a=0.
\]
当 $a=0$ 时，
\[
C\to
\begin{pmatrix}
1&0&0&-2\\
0&1&0&1\\
0&0&1&0\\
0&0&0&0
\end{pmatrix}.
\]
主元列为第 $1,2,3$ 列，对应原矩阵中的 $\alpha_1,\alpha_2,\alpha_3$。由第 $4$ 列关系得
\[
\alpha_4=-2\alpha_1+\alpha_2.
\]

\textbf{【最终答案】}
\[
a=0,\quad \alpha_1,\alpha_2,\alpha_3\text{ 为一个极大线性无关组},
\quad \alpha_4=-2\alpha_1+\alpha_2.
\]

\textbf{【方法总结】} 向量组含参数时，先由秩条件定参数，再读主元列和线性表出。
\end{answerbox}


\section{线性方程组：有解、无穷多解、通解}

\begin{identifybox}[识别]
线性方程组先写增广矩阵，比较 $r(A)$ 与 $r(A|b)$，再确定主元变量和自由变量。
\end{identifybox}
\begin{templatebox}[考场模板]
非齐次通解写成 $X=X_0+k_1\xi_1+k_2\xi_2+\cdots$；齐次方程基础解系个数为 $n-r(A)$。
\end{templatebox}

\subsection{真题融合}

\begin{answerbox}[2020--2021 第一大题第 5 题：填空题]
\textbf{【来源】} 2020--2021 第一大题第 5 题

\textbf{【原题】}
设
\[
A=\begin{pmatrix}1&2&3\\2&4&6\\3&6&9\end{pmatrix},
\]
求三元线性方程组 $AX=0$ 的基础解系中含有几个向量。

\textbf{【速解】}
$\rank(A)=1$，故基础解系含 $3-1=2$ 个向量。

\textbf{【答案】}
\[2\text{ 个向量}\]
\end{answerbox}

\begin{answerbox}[2020--2021 第二大题第 2 题：齐次方程组]
\textbf{【题目分析】} 已知
\[
\begin{cases}
2x_1-3x_2-x_3+x_4=0,\\
-2x_1+2x_2+2x_3-x_4=0,\\
x_2-x_3+2x_4=0.
\end{cases}
\]
写出系数矩阵并求一个基础解系。

\textbf{【详细解答】}
\[
A=\begin{pmatrix}
2&-3&-1&1\\
-2&2&2&-1\\
0&1&-1&2
\end{pmatrix}
\to
\begin{pmatrix}
1&0&-2&0\\
0&1&-1&2\\
0&0&0&0
\end{pmatrix}.
\]
主元变量为 $x_1,x_2$，自由变量为 $x_3,x_4$。

\textbf{【最终答案】}
\[
X=s(2,1,1,0)^T+t(0,-2,0,1)^T.
\]

\textbf{【方法总结】} 齐次方程组基础解系个数为 $n-r(A)$。
\end{answerbox}

\begin{answerbox}[2021--2022 第三大题第 1 题：含参数线性方程组]
\textbf{【题目分析】} 讨论方程组
\[
\begin{cases}
ax_1+x_2+x_3=1,\\
x_1+ax_2+x_3=1,\\
x_1+x_2+ax_3=-2
\end{cases}
\]
何时有解；若有无穷多解，求通解。

\textbf{【详细解答】} 增广矩阵为
\[
\widetilde A=
\begin{pmatrix}
a&1&1&1\\
1&a&1&1\\
1&1&a&-2
\end{pmatrix}.
\]
行化简得
\[
\widetilde A\to
\begin{pmatrix}
1&1&a&-2\\
0&a-1&1-a&3\\
0&0&-(a-1)(a+2)&4+2a
\end{pmatrix}.
\]
当 $a=1$ 时，$r(A)=1<r(A|b)=2$，无解；当 $a=-2$ 时，$r(A)=r(A|b)=2<3$，有无穷多解；当 $a\ne1,-2$ 时，$r(A)=r(A|b)=3$，有唯一解。

当 $a=-2$ 时，继续化简得
\[
\widetilde A\to
\begin{pmatrix}
1&0&-1&-1\\
0&1&-1&-1\\
0&0&0&0
\end{pmatrix},
\]
即
\[
x_1-x_3=-1,\qquad x_2-x_3=-1.
\]
令 $x_3=t$，得
\[
X=\begin{pmatrix}x_1\\x_2\\x_3\end{pmatrix}
=\begin{pmatrix}t-1\\t-1\\t\end{pmatrix}
=t\begin{pmatrix}1\\1\\1\end{pmatrix}
\begin{pmatrix}-1\\-1\\0\end{pmatrix}.
\]

\textbf{【最终答案】}
\[
a\ne1\text{ 时方程组有解};\quad
a=1\text{ 时无解};\quad
a=-2\text{ 时有无穷多解}.
\]
当 $a=-2$ 时，
\[
X=\begin{pmatrix}-1\\-1\\0\end{pmatrix}
t\begin{pmatrix}1\\1\\1\end{pmatrix},\quad t\in\mathbb R.
\]

\textbf{【方法总结】} 参数方程组必须比较 $r(A)$ 与 $r(A|b)$；无穷多解写成“特解 + 齐次基础解系”。
\end{answerbox}

\begin{answerbox}[2022--2023 第一大题第 5 题：填空题]
\textbf{【来源】} 2022--2023 第一大题第 5 题

\textbf{【原题】}
$\alpha_1=(1,4,2)^T,\alpha_2=(2,7,3)^T,\alpha_3=(0,1,3)^T$，求 $x_1\alpha_1+x_2\alpha_2+x_3\alpha_3=0$ 基础解系所含向量个数。

\textbf{【速解】}
列向量矩阵满秩，故齐次方程组只有零解。

\textbf{【答案】}
\[\text{基础解系含 }0\text{ 个向量}\]
\end{answerbox}

\begin{answerbox}[2022--2023 第三大题第 1 题：含参数线性方程组]
\textbf{【题目分析】} 给定方程组
\[
\begin{cases}
2x_1+ax_2-x_3=1,\\
ax_1-x_2+x_3=2,\\
4x_1+5x_2-5x_3=-1
\end{cases}
\]
求：$(1)$ $a$ 为何值时方程组有解；$(2)$ 当方程组有无穷多解时，求通解。

\textbf{【详细解答】} 系数行列式
\[
|A|=(a-1)(5a+4).
\]
当 $a\ne1,-\frac45$ 时，$r(A)=r(A|b)=3$，方程组有唯一解。代入 $a=-\frac45$，有
\[
r(A)\ne r(A|b),
\]
方程组无解。代入 $a=1$，增广矩阵行化简得
\[
(A|b)\to
\begin{pmatrix}
1&0&0&1\\
0&1&-1&-1\\
0&0&0&0
\end{pmatrix}.
\]
故 $r(A)=r(A|b)=2<3$，有无穷多解，且
\[
x_1=1,\qquad x_2-x_3=-1.
\]
令 $x_3=t$，得
\[
X=\begin{pmatrix}1\\-1\\0\end{pmatrix}
+t\begin{pmatrix}0\\1\\1\end{pmatrix}.
\]

\textbf{【最终答案】}
\[
a\ne-\frac45\text{ 时方程组有解};
\quad a=1\text{ 时有无穷多解},
\]
且当 $a=1$ 时
\[
X=\begin{pmatrix}1\\-1\\0\end{pmatrix}
+t\begin{pmatrix}0\\1\\1\end{pmatrix},\quad t\in\mathbb R.
\]

\textbf{【方法总结】} 参数方程组先看行列式，再分别代入退化参数比较 $r(A)$ 与 $r(A|b)$。
\end{answerbox}

\begin{answerbox}[2023--2024 第一大题第 4 题：填空题]
\textbf{【来源】} 2023--2024 第一大题第 4 题

\textbf{【原题】}
设
\[
A=\begin{pmatrix}1&2&-2\\4&t&3\\3&-1&1\end{pmatrix},
\]
$B$ 为 $3$ 阶非零矩阵，且 $AB=0$，求 $t$。

\textbf{【速解】}
$B\ne O$ 且 $AB=0$，故 $|A|=0$。

\textbf{【答案】}
\[t=-3\]
\end{answerbox}

\begin{answerbox}[2023--2024 第三大题第 1 题：参数齐次线性方程组]
\textbf{【题目分析】} 已知 $k$ 是未知常数，线性方程组
\[
\begin{cases}
2x_1-x_2+x_3=0,\\
x_1+kx_2-x_3=0,\\
kx_1+x_2+x_3=0
\end{cases}
\]
什么时候有唯一解、什么时候有无穷多解；若 $k>0$ 且有无穷多解，求其通解。

\textbf{【详细解答】} 系数矩阵行列式为
\[
\begin{vmatrix}
2&-1&1\\
1&k&-1\\
k&1&1
\end{vmatrix}
=-(k-4)(k+1).
\]
故 $k\ne-1,4$ 时，方程组只有零解；$k=-1$ 或 $k=4$ 时，有无穷多解。若 $k>0$ 且有无穷多解，则 $k=4$。此时
\[
\begin{pmatrix}
2&-1&1\\
1&4&-1\\
4&1&1
\end{pmatrix}
\to
\begin{pmatrix}
1&0&\frac13\\
0&1&-\frac13\\
0&0&0
\end{pmatrix}.
\]
令 $x_3=3c$，得
\[
X=c(-1,1,3)^T.
\]

\textbf{【最终答案】}
\[
k\ne-1,4\text{ 时唯一零解};\quad k=-1\text{ 或 }k=4\text{ 时无穷多解}.
\]
若 $k>0$ 且有无穷多解，则
\[
X=c(-1,1,3)^T,\quad c\in\mathbb R.
\]

\textbf{【方法总结】} 齐次方程组“唯一解”就是只有零解；看系数行列式是否为零。
\end{answerbox}

\begin{answerbox}[2024--2025 第三大题第 1 题：非齐次线性方程组]
\textbf{【题目分析】} 已知线性方程组
\[
\begin{cases}
x_1+2x_2+x_3+x_4=1,\\
2x_1+x_2+3x_3+2x_4=3,\\
x_1-x_2+2x_3+x_4=a,\\
3x_1+3x_2+4x_3+3x_4=4
\end{cases}
\]
问 $a$ 为何值时方程组有解；当方程组有解时，求其通解。

\textbf{【详细解答】} 写增广矩阵并行化简，得
\[
a=2
\]
时 $r(A)=r(A|b)=2<4$，方程组有无穷多解；$a\ne2$ 时 $r(A)<r(A|b)$，无解。$a=2$ 时，主元变量为 $x_1,x_2$，自由变量为 $x_3,x_4$，通解为
\[
X=\left(\frac53,-\frac13,0,0\right)^T
+c_1(-5,1,3,0)^T+c_2(-1,0,0,1)^T.
\]

\textbf{【最终答案】}
\[
a=2\text{ 时有解，}a\ne2\text{ 时无解};
\]
当 $a=2$ 时，
\[
X=
\left(\frac53,-\frac13,0,0\right)^T
+c_1(-5,1,3,0)^T
+c_2(-1,0,0,1)^T.
\]

\textbf{【方法总结】} 非齐次方程组通解必须写成“特解 + 齐次方程组基础解系线性组合”。
\end{answerbox}


\section{特征值、特征向量、相似对角化、实对称矩阵}

\begin{identifybox}[识别]
先写特征方程，再分别求 $(A-\lambda E)X=0$ 的特征子空间；实对称矩阵可正交对角化。
\end{identifybox}
\begin{templatebox}[考场模板]
若 $P=(\xi_1,\ldots,\xi_n)$，则 $P^{-1}AP=\Lambda$，$P$ 的列向量顺序必须与 $\Lambda$ 对角元一致；实对称矩阵写 $Q^TAQ=\Lambda$。
\end{templatebox}

\subsection{真题融合}

\begin{answerbox}[2020--2021 第一大题第 4 题：填空题]
\textbf{【来源】} 2020--2021 第一大题第 4 题

\textbf{【原题】}
已知实对称矩阵
\[
\begin{pmatrix}a&1&1\\1&a&1\\1&1&a\end{pmatrix}
\]
与对角矩阵 $\diag(3,0,0)$ 相似，求 $a$。

\textbf{【速解】}
特征值为 $a+2,a-1,a-1$，与 $3,0,0$ 对应，得 $a=1$。

\textbf{【答案】}
\[a=1\]
\end{answerbox}

\begin{answerbox}[2020--2021 第三大题第 2 题：实对称矩阵]
\textbf{【题目分析】} 设 $A$ 为 $3$ 阶实对称矩阵，$\rank(A)=2$，且
\[
A\begin{pmatrix}1&-2\\0&0\\1&2\end{pmatrix}
=
\begin{pmatrix}1&2\\0&0\\1&-2\end{pmatrix}.
\]
求 $A$ 的特征值、特征向量，并求 $A$。

\textbf{【详细解答】} 由题给关系得到两个特征向量：
\[
(1,0,1)^T\mapsto (1,0,1)^T,\qquad
(-2,0,2)^T\mapsto (2,0,-2)^T.
\]
故对应特征值为 $1,-1$。又 $\rank(A)=2$，第三个特征值为 $0$。取单位正交特征向量
\[
\gamma_1=\frac1{\sqrt2}(1,0,1)^T,\quad
\gamma_2=\frac1{\sqrt2}(-1,0,1)^T,\quad
\gamma_3=(0,1,0)^T.
\]
令 $Q=(\gamma_1,\gamma_2,\gamma_3)$，则
\[
Q^TAQ=\diag(1,-1,0),\qquad
A=Q\diag(1,-1,0)Q^T.
\]

\textbf{【最终答案】} 特征值为 $1,-1,0$，
\[
A=\begin{pmatrix}0&0&1\\0&0&0\\1&0&0\end{pmatrix}.
\]

\textbf{【方法总结】} 实对称矩阵用正交谱分解，秩等于非零特征值个数。
\end{answerbox}

\begin{answerbox}[2021--2022 第一大题第 4 题：填空题]
\textbf{【来源】} 2021--2022 第一大题第 4 题

\textbf{【原题】}
$A=\begin{pmatrix}2&1&a\\1&-1&2\\3&0&1\end{pmatrix}$，$B$ 为 $3$ 阶非零方阵，$AB=0$，求 $a$。

\textbf{【速解】}
$AB=0$ 且 $B\ne O$，故 $|A|=0$。

\textbf{【答案】}
\[a=-1\]
\end{answerbox}

\begin{answerbox}[2021--2022 第三大题第 2 题：实对称矩阵]
\textbf{【题目分析】} $A$ 为 $3$ 阶实对称矩阵，特征值为
\[
\lambda_1=-1,\qquad \lambda_2=\lambda_3=1,
\]
且 $\xi_1=(0,1,1)^T$ 是 $\lambda_1=-1$ 的特征向量，求 $A$。

\textbf{【详细解答】} 实对称矩阵不同特征值对应特征向量正交。属于特征值 $1$ 的特征向量 $\xi=(x_1,x_2,x_3)^T$ 满足
\[
(\xi,\xi_1)=x_2+x_3=0.
\]
可取属于 $1$ 的两个线性无关特征向量
\[
\xi_2=(1,0,0)^T,\qquad \xi_3=(0,-1,1)^T.
\]
单位化得
\[
\gamma_1=\left(0,\frac{\sqrt2}{2},\frac{\sqrt2}{2}\right)^T,\quad
\gamma_2=(1,0,0)^T,\quad
\gamma_3=\left(0,-\frac{\sqrt2}{2},\frac{\sqrt2}{2}\right)^T.
\]
令 $P=(\gamma_1,\gamma_2,\gamma_3)$，则 $P$ 为正交矩阵，且
\[
P^{-1}AP=P^TAP=\diag(-1,1,1).
\]
于是
\[
A=P\diag(-1,1,1)P^T
=\begin{pmatrix}1&0&0\\0&0&-1\\0&-1&0\end{pmatrix}.
\]

\textbf{【最终答案】}
\[
A=\begin{pmatrix}1&0&0\\0&0&-1\\0&-1&0\end{pmatrix}.
\]

\textbf{【方法总结】} 实对称矩阵重特征值子空间内补足正交基，再用 $A=Q\Lambda Q^T$。
\end{answerbox}

\begin{answerbox}[2022--2023 第三大题第 2 题：特征值、对角化与矩阵幂]
\textbf{【题目分析】} 已知
\[
A=\begin{pmatrix}1&1&2\\0&a&2\\1&-1&0\end{pmatrix},
\]
且 $\lambda=0$ 是 $A$ 的特征值。求：$(1)$ $a$ 的值；$(2)$ $A$ 的所有特征值；$(3)$ $A^{100}$。

\textbf{【详细解答】} 因 $\lambda=0$ 是特征值，故 $|A|=0$：
\[
|A|=\begin{vmatrix}1&1&2\\0&a&2\\1&-1&0\end{vmatrix}=-2a+4=0,
\]
得 $a=2$。此时
\[
|\lambda I-A|=\lambda(\lambda-1)(\lambda-2),
\]
故特征值为 $0,1,2$。分别解 $(\lambda I-A)X=0$，可取
\[
\xi_1=(-1,-1,1)^T,\quad
\xi_2=(-1,-2,1)^T,\quad
\xi_3=(1,1,0)^T.
\]
令 $P=(\xi_1,\xi_2,\xi_3)$，则
\[
P^{-1}AP=\Lambda=\diag(0,1,2),
\quad
P^{-1}=\begin{pmatrix}-1&1&1\\1&-1&0\\1&0&1\end{pmatrix}.
\]
于是 $A^{100}=P\Lambda^{100}P^{-1}$。

\textbf{【最终答案】}
\[
a=2,\qquad \lambda=0,1,2,
\]
\[
A^{100}=
\begin{pmatrix}
2^{100}-1&1&2^{100}\\
2^{100}-2&2&2^{100}\\
1&-1&0
\end{pmatrix}.
\]

\textbf{【方法总结】} 矩阵幂题先对角化，$P$ 的列向量顺序必须与 $\Lambda$ 对角元顺序一致。
\end{answerbox}

\begin{answerbox}[2023--2024 第一大题第 3 题：填空题]
\textbf{【来源】} 2023--2024 第一大题第 3 题

\textbf{【原题】}
已知 $3$ 阶矩阵 $A$ 的特征值为 $1,2,3$，$I$ 是 $3$ 阶单位矩阵，求 $|4I-A|$。

\textbf{【速解】}
$4I-A$ 的特征值为 $4-\lambda_i$。

\textbf{【答案】}
\[|4I-A|=6\]
\end{answerbox}

\begin{answerbox}[2023--2024 第一大题第 6 题：填空题]
\textbf{【来源】} 2023--2024 第一大题第 6 题

\textbf{【原题】}
设
\[
B=\begin{pmatrix}1&2&2\\0&2&3\\0&1&\lambda\end{pmatrix},
\]
已知二次型 $f(x)=x^TBx$ 是正定的，求 $\lambda$ 的取值范围。

\textbf{【速解】}
二次型矩阵取对称部分，按正定主子式判别。

\textbf{【答案】}
\[\lambda\in(2,+\infty)\]
\end{answerbox}

\begin{answerbox}[2023--2024 第三大题第 2 题：特征向量参数与对角化]
\textbf{【题目分析】} 已知
\[
\xi=(1,1,-1)^T
\]
是矩阵
\[
A=\begin{pmatrix}2&-1&2\\5&a&3\\-1&b&-2\end{pmatrix}
\]
的一个特征向量。$(1)$ 确定参数 $a,b$ 及 $\xi$ 对应的特征值；$(2)$ 问 $A$ 能否相似于对角阵，并说明理由。

\textbf{【详细解答】} 设 $A\xi=\lambda\xi$，则
\[
A\xi=\begin{pmatrix}-1\\5+a-3\\-1+b+2\end{pmatrix}
=\begin{pmatrix}-1\\a+2\\b+1\end{pmatrix}
=\lambda(1,1,-1)^T.
\]
比较分量得
\[
\lambda=-1,\qquad a=-3,\qquad b=0.
\]
代入后
\[
A=\begin{pmatrix}2&-1&2\\5&-3&3\\-1&0&-2\end{pmatrix}.
\]
特征多项式为
\[
|\lambda I-A|=(\lambda+1)^3.
\]
全部特征值均为 $-1$。对 $\lambda=-1$，解 $(-I-A)X=0$，其系数矩阵行化简为
\[
-I-A\to
\begin{pmatrix}
1&0&1\\
0&1&1\\
0&0&0
\end{pmatrix},
\]
故特征子空间维数为 $1$。

\textbf{【最终答案】}
\[
\lambda=-1,\quad a=-3,\quad b=0.
\]
$A$ 只有 $1$ 个线性无关特征向量，不能相似于对角阵。

\textbf{【方法总结】} 对角化要比较特征子空间维数与矩阵阶数，不能只看特征值。
\end{answerbox}

\begin{answerbox}[2024--2025 第一大题第 6 题：填空题]
\textbf{【来源】} 2024--2025 第一大题第 6 题

\textbf{【原题】}
设二次型
\[
f(x_1,x_2,x_3)=5x_1^2+5x_2^2+kx_3^2-4x_1x_2+4x_1x_3-8x_2x_3
\]
为正定二次型，求实数 $k$ 的取值范围。

\textbf{【速解】}
写二次型矩阵并用顺序主子式判别。

\textbf{【答案】}
\[k>\frac{68}{21}\]
\end{answerbox}

\begin{answerbox}[2024--2025 第三大题第 3 题：特征多项式与向量方程]
\textbf{【题目分析】} 设
\[
A=\begin{pmatrix}0&-1&2\\-1&0&2\\-1&-1&k\end{pmatrix}.
\]
已知 $1$ 是 $A$ 的特征多项式的重根。求 $k$ 的值；求所有满足
\[
A\alpha=\alpha+\beta,\qquad A^2\alpha=\alpha+2\beta
\]
的非零列向量 $\alpha,\beta$。

\textbf{【详细解答】} 由特征多项式中 $1$ 为重根，得
\[
k=3.
\]
此时
\[
(A-I)X=0
\]
的基础解系可取
\[
\xi_1=(-1,1,0)^T,\qquad \xi_2=(2,0,1)^T.
\]
由
\[
A^2\alpha=A(\alpha+\beta)=A\alpha+A\beta=\alpha+\beta+A\beta
\]
和 $A^2\alpha=\alpha+2\beta$ 得
\[
A\beta=\beta,
\]
故
\[
\beta=c\bigl((-1,1,0)^T+(2,0,1)^T\bigr),\quad c\ne0,
\]
并由 $(A-I)\alpha=\beta$ 解得
\[
\alpha=(-c,0,0)^T+s(-1,1,0)^T+t(2,0,1)^T.
\]

\textbf{【最终答案】}
\[
k=3,
\]
\[
\beta=c(1,1,1)^T,\quad c\ne0,
\]
\[
\alpha=(-c,0,0)^T+s(-1,1,0)^T+t(2,0,1)^T,\quad s,t\in\mathbb R.
\]

\textbf{【方法总结】} 由两式相减先推出 $\beta$ 属于特征值 $1$ 的特征子空间，再解非齐次方程 $(A-I)\alpha=\beta$。
\end{answerbox}


\section{二次型：正定、标准形、正交变换}

\begin{identifybox}[识别]
二次型先写对称矩阵，再选择配方法、合同变换或正交变换；正定题看顺序主子式或特征值。
\end{identifybox}
\begin{templatebox}[考场模板]
写清标准形、规范形、正惯性指数、负惯性指数和秩；正交变换题必须构造单位正交特征向量矩阵 $Q$。
\end{templatebox}

\subsection{真题融合}

\begin{answerbox}[2020--2021 第一大题第 6 题：填空题]
\textbf{【来源】} 2020--2021 第一大题第 6 题

\textbf{【原题】}
设
\[
A=\begin{pmatrix}1&t&2\\t&4&t\\2&t&5\end{pmatrix}
\]
为正定矩阵，求 $t$ 的取值范围。

\textbf{【速解】}
顺序主子式均大于 $0$，得 $-\sqrt2<t<\sqrt2$。

\textbf{【答案】}
\[-\sqrt2<t<\sqrt2\]
\end{answerbox}

\begin{answerbox}[2020--2021 第三大题第 3 题：二次型正交化]
\textbf{【题目分析】} 已知
\[
f=x_1^2+4x_2^2+x_3^2-4x_1x_2-8x_1x_3-4x_2x_3.
\]
写矩阵 $A$，并求正交矩阵 $Q$ 使 $x=Qy$ 后化为标准形。

\textbf{【详细解答】}
\[
A=\begin{pmatrix}
1&-2&-4\\
-2&4&-2\\
-4&-2&1
\end{pmatrix}.
\]
求 $|\lambda I-A|=0$，再对每个特征值求单位正交特征向量，列成 $Q$，则
\[
Q^TAQ=\Lambda.
\]
计算得
\[
|\lambda I-A|=(\lambda-5)^2(\lambda+4).
\]
取单位正交特征向量
\[
\gamma_1=\frac1{\sqrt2}(1,0,-1)^T,\quad
\gamma_2=\frac1{3\sqrt2}(1,-4,1)^T,\quad
\gamma_3=\frac13(2,1,2)^T.
\]
令 $Q=(\gamma_1,\gamma_2,\gamma_3)$，则
\[
Q^TAQ=\diag(5,5,-4).
\]

\textbf{【最终答案】}
\[
Q=\begin{pmatrix}
\frac1{\sqrt2}&\frac1{3\sqrt2}&\frac23\\
0&-\frac4{3\sqrt2}&\frac13\\
-\frac1{\sqrt2}&\frac1{3\sqrt2}&\frac23
\end{pmatrix},
\qquad
f=5y_1^2+5y_2^2-4y_3^2.
\]

\textbf{【方法总结】} 二次型交叉项 $2a_{ij}x_ix_j$，写矩阵时交叉项系数要除以 $2$。
\end{answerbox}

\begin{answerbox}[2021--2022 第一大题第 6 题：填空题]
\textbf{【来源】} 2021--2022 第一大题第 6 题

\textbf{【原题】}
$A=\begin{pmatrix}2&1&0\\1&1&t\\0&t&1\end{pmatrix}$ 为正定矩阵，求 $t$ 的取值范围。

\textbf{【速解】}
正定矩阵按顺序主子式判别。

\textbf{【答案】}
\[-\frac{\sqrt2}{2}<t<\frac{\sqrt2}{2}\]
\end{answerbox}

\begin{answerbox}[2021--2022 第三大题第 3 题：二次型正交化]
\textbf{【题目分析】} 已知
\[
f=2x_1^2+5x_2^2+5x_3^2+4x_1x_2-4x_1x_3-8x_2x_3.
\]
写出二次型矩阵，并求正交变换将其化为标准形。

\textbf{【详细解答】} 二次型矩阵为
\[
A=\begin{pmatrix}2&2&-2\\2&5&-4\\-2&-4&5\end{pmatrix}.
\]
特征多项式
\[
|\lambda I-A|
=\begin{vmatrix}
\lambda-2&-2&2\\
-2&\lambda-5&4\\
2&4&\lambda-5
\end{vmatrix}
=(\lambda-1)^2(\lambda-10).
\]
故特征值为 $1,1,10$。当 $\lambda=1$ 时，可取基础解系
\[
\xi_1=(2,0,1)^T,\qquad \xi_2=(-2,1,0)^T.
\]
施密特正交化并单位化，取
\[
\gamma_1=\frac1{\sqrt5}(2,0,1)^T,\qquad
\gamma_2=\frac{\sqrt5}{3}\left(-\frac25,1,\frac45\right)^T.
\]
当 $\lambda=10$ 时，可取单位特征向量
\[
\gamma_3=\left(-\frac13,-\frac23,\frac23\right)^T.
\]
令
\[
Q=(\gamma_1,\gamma_2,\gamma_3)
=\begin{pmatrix}
\frac{2\sqrt5}{5}&-\frac{2\sqrt5}{15}&-\frac13\\
0&\frac{\sqrt5}{3}&-\frac23\\
\frac{\sqrt5}{5}&\frac{4\sqrt5}{15}&\frac23
\end{pmatrix},
\]
则
\[
Q^TAQ=\diag(1,1,10).
\]

\textbf{【最终答案】}
\[
x=Qy,\qquad f=y_1^2+y_2^2+10y_3^2.
\]

\textbf{【方法总结】} 二重特征值下先取基础解系，再施密特正交化、单位化。
\end{answerbox}

\begin{answerbox}[2022--2023 第一大题第 6 题：填空题]
\textbf{【来源】} 2022--2023 第一大题第 6 题

\textbf{【原题】}
$A=\begin{pmatrix}1&t&-1\\t&4&2\\-1&2&4\end{pmatrix}$ 为正定矩阵，求 $t$ 范围。

\textbf{【速解】}
顺序主子式均大于 $0$。

\textbf{【答案】}
\[-2<t<1\]
\end{answerbox}

\begin{answerbox}[2022--2023 第三大题第 3 题：二次型正交变换]
\textbf{【题目分析】} 已知二次型
\[
f(x_1,x_2,x_3)=2x_1^2+3x_2^2+3x_3^2+2ax_2x_3\quad(a>0)
\]
通过正交变换可化为标准形
\[
f=y_1^2+2y_2^2+5y_3^2.
\]
求二次型矩阵、参数 $a$ 及正交矩阵 $Q$。

\textbf{【详细解答】} 二次型矩阵为
\[
A=\begin{pmatrix}2&0&0\\0&3&a\\0&a&3\end{pmatrix}.
\]
由标准形知 $A$ 的特征值为 $1,2,5$，故
\[
|A|=2(9-a^2)=1\cdot2\cdot5=10.
\]
又 $a>0$，得 $a=2$。对应特征值 $1,2,5$ 的单位正交特征向量可取
\[
\gamma_1=\left(0,\frac{\sqrt2}{2},-\frac{\sqrt2}{2}\right)^T,\quad
\gamma_2=(1,0,0)^T,\quad
\gamma_3=\left(0,\frac{\sqrt2}{2},\frac{\sqrt2}{2}\right)^T.
\]
令 $Q=(\gamma_1,\gamma_2,\gamma_3)$，则
\[
Q=\begin{pmatrix}
0&1&0\\
\frac{\sqrt2}{2}&0&\frac{\sqrt2}{2}\\
-\frac{\sqrt2}{2}&0&\frac{\sqrt2}{2}
\end{pmatrix},
\quad Q^TAQ=\diag(1,2,5).
\]

\textbf{【最终答案】}
\[
a=2,\qquad
Q=\begin{pmatrix}
0&1&0\\
\frac{\sqrt2}{2}&0&\frac{\sqrt2}{2}\\
-\frac{\sqrt2}{2}&0&\frac{\sqrt2}{2}
\end{pmatrix},
\quad f=y_1^2+2y_2^2+5y_3^2.
\]

\textbf{【方法总结】} 已知标准形时，标准形系数就是实对称矩阵的特征值。
\end{answerbox}

\begin{answerbox}[2023--2024 第三大题第 3 题：二次型正交化]
\textbf{【题目分析】} 已知二次型
\[
f(x_1,x_2,x_3)=4x_2^2-3x_3^2+4x_1x_2-4x_1x_3+8x_2x_3.
\]
$(1)$ 写出二次型矩阵；$(2)$ 用正交变换化为标准形，并写出相应正交矩阵。

\textbf{【详细解答】} 二次型矩阵为
\[
A=\begin{pmatrix}
0&2&-2\\
2&4&4\\
-2&4&-3
\end{pmatrix}.
\]
特征多项式为
\[
|\lambda I-A|=(\lambda+6)(\lambda-1)(\lambda-6),
\]
故特征值为 $-6,1,6$。对应特征向量可取
\[
v_1=(1,-1,2)^T,\quad v_2=(-2,0,1)^T,\quad v_3=(1,5,2)^T.
\]
因 $A$ 为实对称矩阵，不同特征值对应特征向量正交，单位化得
\[
Q=\begin{pmatrix}
\frac1{\sqrt6}&-\frac2{\sqrt5}&\frac1{\sqrt{30}}\\
-\frac1{\sqrt6}&0&\frac5{\sqrt{30}}\\
\frac2{\sqrt6}&\frac1{\sqrt5}&\frac2{\sqrt{30}}
\end{pmatrix}.
\]
则
\[
Q^TAQ=\diag(-6,1,6).
\]

\textbf{【最终答案】}
\[
A=\begin{pmatrix}0&2&-2\\2&4&4\\-2&4&-3\end{pmatrix},
\quad
f=-6y_1^2+y_2^2+6y_3^2.
\]

\textbf{【方法总结】} 实对称二次型正交化：矩阵、特征值、单位正交特征向量、$Q^TAQ$、标准形都要写。
\end{answerbox}

\begin{answerbox}[2024--2025 第一大题第 3 题：填空题]
\textbf{【来源】} 2024--2025 第一大题第 3 题

\textbf{【原题】}
设
\[
A=\begin{pmatrix}1&-1&a\\2&a&-4\\3&1&0\end{pmatrix}.
\]
若三元齐次线性方程组 $Ax=0$ 有非零解，求 $a$。

\textbf{【速解】}
有非零解等价于 $|A|=0$。

\textbf{【答案】}
\[a=-2\text{ 或 }a=\frac83\]
\end{answerbox}

\begin{answerbox}[2024--2025 第三大题第 2 题：二次型正交化]
\textbf{【题目分析】} 已知二次型
\[
f=5x_1^2+5x_2^2+2x_3^2-8x_1x_2+4x_1x_3-4x_2x_3,
\]
写出二次型矩阵 $A$，求正交矩阵 $Q$，使 $x=Qy$ 将二次型化为标准形，并写出该标准形。

\textbf{【详细解答】} 二次型矩阵为
\[
A=\begin{pmatrix}5&-4&2\\-4&5&-2\\2&-2&2\end{pmatrix}.
\]
特征值为 $1,1,10$，可取
\[
Q=\begin{pmatrix}
\frac{\sqrt2}{2}&-\frac{\sqrt2}{6}&\frac23\\
\frac{\sqrt2}{2}&\frac{\sqrt2}{6}&-\frac23\\
0&\frac{2\sqrt2}{3}&\frac13
\end{pmatrix},
\quad Q^TAQ=\diag(1,1,10).
\]
标准形为 $y_1^2+y_2^2+10y_3^2$。

\textbf{【最终答案】}
\[
Q^TAQ=\diag(1,1,10),\qquad f=y_1^2+y_2^2+10y_3^2.
\]

\textbf{【方法总结】} 实对称二次型必须写二次型矩阵、正交矩阵和标准形。
\end{answerbox}


\section{证明题：固定套路}

\begin{identifybox}[识别]
证明题先判断要证明线性无关、可逆、相似、正交还是秩结论，再选对应定义和常用恒等式。
\end{identifybox}
\begin{templatebox}[考场模板]
线性无关：设线性组合为零；可逆：证行列式非零或分解为可逆矩阵乘积；伴随矩阵：用 $AA^*=A^*A=|A|E$。
\end{templatebox}

\subsection{真题融合}

\begin{answerbox}[2020--2021 第四大题第 1 题：证明题]
\textbf{【原题】} 设 $A$ 与 $B$ 均为 $3$ 阶方阵且 $A$ 与 $B$ 相似。证明：
\[
A^2+2A+3I_3\text{ 与 }B^2+2B+3I_3\text{ 相似}.
\]

\textbf{【证明思路】} 若 $A,B$ 相似，则存在可逆矩阵 $P$ 使 $B=P^{-1}AP$，从而 $B$ 的多项式可写为 $A$ 的同一多项式的相似变换。

\textbf{【详细证明】} 设 $B=P^{-1}AP$，其中 $P$ 可逆，则
\[
B^2=P^{-1}A^2P,\qquad I=P^{-1}IP.
\]
于是
\[
B^2+2B+3I=P^{-1}(A^2+2A+3I)P,
\]
故 $B^2+2B+3I$ 与 $A^2+2A+3I$ 相似。

\textbf{【结论】} 矩阵多项式保持相似关系。
\end{answerbox}

\begin{answerbox}[2020--2021 第四大题第 2 题：证明题]
\textbf{【原题】} 设 $A$ 为 $3$ 阶方阵，$\alpha_1,\alpha_2$ 为 $A$ 的分别属于 $1,-1$ 的特征值的特征向量，向量 $\alpha_3$ 满足
\[
A\alpha_3=\alpha_1+\alpha_3.
\]
证明：$\alpha_1,\alpha_2,\alpha_3$ 线性无关。

\textbf{【证明思路】} 若 $\alpha_1,\alpha_2$ 分别属于不同特征值，且 $A\alpha_3=\alpha_1+\alpha_3$，用线性组合为零并作用 $A$，再与原式联立证明系数全为零。

\textbf{【详细证明】} 设
\[
c_1\alpha_1+c_2\alpha_2+c_3\alpha_3=0.
\]
对等式作用 $A$，结合 $A\alpha_1=\alpha_1,\ A\alpha_2=-\alpha_2,\ A\alpha_3=\alpha_1+\alpha_3$，与原式联立可得 $c_1=c_2=c_3=0$。

\textbf{【结论】} $\alpha_1,\alpha_2,\alpha_3$ 线性无关。
\end{answerbox}

\begin{answerbox}[2021--2022 第四大题第 1 题：证明题]
\textbf{【原题】} 设 $A$ 为 $4\times3$ 型的实矩阵，且 $A^TA$ 的对角线元素之和为零，其中 $A^T$ 表示 $A$ 的转置。证明：$A$ 为 $4\times3$ 型的零矩阵。

\textbf{【证明思路】} 设 $A$ 为 $4\times3$ 实矩阵，若 $A^TA$ 对角线元素之和为 $0$，证明 $A=O$。核心是把 $A^TA$ 的对角线和写成 $A$ 所有元素平方和。

\textbf{【详细证明】}
设
\[
A=(a_{ij})_{4\times3},\qquad C=A^TA.
\]
则对 $1\le j\le3$，
\[
C(j,j)=\sum_{i=1}^4 a_{ij}^2.
\]
由题意，
\[
\sum_{j=1}^3 C(j,j)=\sum_{j=1}^3\sum_{i=1}^4 a_{ij}^2=0.
\]
各项均为非负数，故 $a_{ij}=0$，从而 $A=O$。

\textbf{【结论】} $A$ 为 $4\times3$ 零矩阵。
\end{answerbox}

\begin{answerbox}[2021--2022 第四大题第 2 题：证明题]
\textbf{【原题】} 设某个 $4$ 元非齐次线性方程组的系数矩阵的秩为 $2$，且该方程组有无穷多个解。证明：存在该方程组的 $3$ 个线性无关的解向量 $\gamma_1,\gamma_2,\gamma_3$，使得该方程组的任意解向量 $\gamma$ 都可以由 $\gamma_1,\gamma_2,\gamma_3$ 线性表出。

\textbf{【证明思路】} 四元非齐次线性方程组 $Ax=b\ (b\ne0)$ 满足
\[
r(A)=r(A|b)=2<4.
\]
证明存在 $3$ 个线性无关解向量，使任一解可由其线性表出。用“一个特解 + 两个基础解系向量”构造。

\textbf{【详细证明】} 由 $r(A)=2$，齐次方程组 $Ax=0$ 的基础解系含 $4-2=2$ 个向量，设为 $\eta_1,\eta_2$。取 $Ax=b$ 的一个特解 $\eta$。令
\[
\gamma_1=\eta,\qquad
\gamma_2=\eta+\eta_1,\qquad
\gamma_3=\eta+\eta_2.
\]
则 $\gamma_1,\gamma_2,\gamma_3$ 均为 $Ax=b$ 的解。任一解向量可写为
\[
\gamma=\eta+k_1\eta_1+k_2\eta_2
=\gamma_1+k_1(\gamma_2-\gamma_1)+k_2(\gamma_3-\gamma_1),
\]
故任一解可由 $\gamma_1,\gamma_2,\gamma_3$ 线性表出。

再证线性无关。设
\[
x_1\gamma_1+x_2\gamma_2+x_3\gamma_3=0.
\]
代入 $\gamma_i$ 得
\[
(x_1+x_2+x_3)\eta+x_2\eta_1+x_3\eta_2=0. \tag{1}
\]
两边左乘 $A$：
\[
(x_1+x_2+x_3)b=0.
\]
由于 $b\ne0$，故
\[
x_1+x_2+x_3=0. \tag{2}
\]
由 (1)(2) 得 $x_2\eta_1+x_3\eta_2=0$。因 $\eta_1,\eta_2$ 线性无关，故 $x_2=x_3=0$，再由 (2) 得 $x_1=0$。

\textbf{【结论】} 存在 $3$ 个线性无关解向量 $\gamma_1,\gamma_2,\gamma_3$，且任一解均可由它们线性表出。
\end{answerbox}

\begin{answerbox}[2022--2023 第四大题第 1 题：证明题]
\textbf{【原题】} 已知 $\alpha_1,\alpha_2,\alpha_3$ 均为 $n$ 维的非零向量且它们两两正交。证明：$\alpha_1,\alpha_2,\alpha_3$ 线性无关。

\textbf{【证明思路】} 已知 $\alpha_1,\alpha_2,\alpha_3$ 均为 $n$ 维非零向量且两两正交。用线性组合为零，再分别与每个向量作内积。

\textbf{【详细证明】} 设
\[
x_1\alpha_1+x_2\alpha_2+x_3\alpha_3=0.
\]
两边分别与 $\alpha_i$ 作内积，因两两正交，得
\[
x_i\langle\alpha_i,\alpha_i\rangle=0,\quad i=1,2,3.
\]
由于 $\alpha_i$ 非零，故 $\langle\alpha_i,\alpha_i\rangle>0$，于是
\[
x_1=x_2=x_3=0.
\]

\textbf{【结论】} $\alpha_1,\alpha_2,\alpha_3$ 线性无关。
\end{answerbox}

\begin{answerbox}[2022--2023 第四大题第 2 题：证明题]
\textbf{【原题】} 设 $A$ 为 $n$ 阶方阵，$\lambda_1$ 与 $\lambda_2$ 为 $A$ 的两个不同的特征值，而 $\alpha_1,\alpha_2$ 分别属于 $\lambda_1$ 和 $\lambda_2$ 的特征向量。证明：$2\alpha_1+3\alpha_2$ 不是 $A$ 的特征向量。

\textbf{【证明思路】} 设 $A$ 为 $n$ 阶方阵，$\lambda_1,\lambda_2$ 为 $A$ 的两个不同特征值，$\alpha_1,\alpha_2$ 分别属于 $\lambda_1,\lambda_2$。用反证法证明 $2\alpha_1+3\alpha_2$ 不是特征向量。

\textbf{【详细证明】} 若 $2\alpha_1+3\alpha_2$ 是 $A$ 的特征向量，对应特征值 $\lambda$，则
\[
A(2\alpha_1+3\alpha_2)=\lambda(2\alpha_1+3\alpha_2)
=2\lambda\alpha_1+3\lambda\alpha_2. \tag{1}
\]
另一方面，由 $\alpha_1,\alpha_2$ 的特征向量性质，
\[
A(2\alpha_1+3\alpha_2)=2\lambda_1\alpha_1+3\lambda_2\alpha_2. \tag{2}
\]
由 (1)(2) 得
\[
2(\lambda_1-\lambda)\alpha_1+3(\lambda_2-\lambda)\alpha_2=0.
\]
不同特征值对应的特征向量线性无关，故
\[
\lambda_1=\lambda,\qquad \lambda_2=\lambda,
\]
从而 $\lambda_1=\lambda_2$，与题设矛盾。

\textbf{【结论】} $2\alpha_1+3\alpha_2$ 不是 $A$ 的特征向量。
\end{answerbox}

\begin{answerbox}[2023--2024 第四大题第 1 题：证明题]
\textbf{【原题】} 设 $\alpha$ 为 $n$ 维单位列向量，$I$ 为 $n$ 阶单位矩阵。证明：$I-\alpha\alpha^T$ 不可逆。

\textbf{【证明思路】} 设 $\alpha$ 为 $n$ 维单位列向量，证明 $I-\alpha\alpha^T$ 不可逆。直接构造非零向量 $\alpha$ 被该矩阵映为零。

\textbf{【详细证明】} 因 $\alpha$ 为单位列向量，
\[
\alpha^T\alpha=1.
\]
于是
\[
(I-\alpha\alpha^T)\alpha
=\alpha-\alpha(\alpha^T\alpha)
=\alpha-\alpha=0.
\]
又 $\alpha\ne0$，故齐次方程
\[
(I-\alpha\alpha^T)x=0
\]
有非零解。

\textbf{【结论】} $I-\alpha\alpha^T$ 不可逆。
\end{answerbox}

\begin{answerbox}[2023--2024 第四大题第 2 题：证明题]
\textbf{【原题】} 设 $n$ 阶方阵 $A$ 满足
\[
A^2+2A-3I=0,
\]
其中 $I$ 是 $n$ 阶单位矩阵。证明：$(1)$ $A$ 的特征值只可能是 $1$ 或者 $-3$；$(2)$ $A$ 一定可以相似对角化。

\textbf{【证明思路】} 设 $A^2+2A-3I=0$。先令特征向量代入矩阵多项式，得到特征值只可能为 $1$ 或 $-3$；再按参考答案用特征子空间维数证明可对角化。

\textbf{【详细证明】} 若 $Av=\lambda v,\ v\ne0$，则
\[
0=(A^2+2A-3I)v=(\lambda^2+2\lambda-3)v.
\]
故
\[
\lambda^2+2\lambda-3=0,
\]
即 $\lambda=1$ 或 $\lambda=-3$。

由题设
\[
(A-I)(A+3I)=O.
\]
若只有一个特征值，则 $A=I$ 或 $A=-3I$，显然可对角化。若两个特征值均存在，则属于 $1$ 的线性无关特征向量个数为 $n-r(A-I)$，属于 $-3$ 的线性无关特征向量个数为 $n-r(A+3I)$。由
\[
(A-I)(A+3I)=O
\]
得
\[
r(A-I)+r(A+3I)\le n.
\]
于是
\[
[n-r(A-I)]+[n-r(A+3I)]\ge n.
\]
故 $A$ 有 $n$ 个线性无关特征向量。

\textbf{【结论】} $A$ 的特征值只可能是 $1$ 或 $-3$，且 $A$ 一定可以相似对角化。
\end{answerbox}

\begin{answerbox}[2024--2025 第四大题第 1 题：证明题]
\textbf{【原题】} 设向量组 $\alpha_1,\alpha_2,\alpha_3$ 线性无关，证明向量组
\[
\beta_1=\alpha_1+\alpha_2,\quad
\beta_2=\alpha_2+\alpha_3,\quad
\beta_3=\alpha_3+\alpha_1
\]
也线性无关。

\textbf{【证明思路】} 设 $\alpha_1,\alpha_2,\alpha_3$ 线性无关，证明
\[
\beta_1=\alpha_1+\alpha_2,\quad \beta_2=\alpha_2+\alpha_3,\quad \beta_3=\alpha_3+\alpha_1
\]
也线性无关。写成矩阵乘法并证明系数矩阵可逆。

\textbf{【详细证明】} 记
\[
(\beta_1,\beta_2,\beta_3)=(\alpha_1,\alpha_2,\alpha_3)
\begin{pmatrix}1&0&1\\1&1&0\\0&1&1\end{pmatrix}.
\]
右侧系数矩阵行列式为
\[
\begin{vmatrix}1&0&1\\1&1&0\\0&1&1\end{vmatrix}=2\ne0,
\]
故该系数矩阵可逆。由于 $\alpha_1,\alpha_2,\alpha_3$ 线性无关，右乘可逆矩阵不改变列向量组的秩，所以 $\beta_1,\beta_2,\beta_3$ 线性无关。

\textbf{【结论】} $\beta_1,\beta_2,\beta_3$ 线性无关。
\end{answerbox}

\begin{answerbox}[2024--2025 第四大题第 2 题：证明题]
\textbf{【原题】} 已知 $3$ 阶实矩阵 $A$ 的特征值为 $\lambda_1\ge\lambda_2\ge\lambda_3$。证明：$(1)$ 当 $A$ 是对角矩阵时，对于任意 $3$ 维单位向量 $x$，有
\[
\lambda_3\le x^TAx\le \lambda_1;
\]
$(2)$ 当 $A$ 是实对称矩阵时，对于任意 $3$ 维单位向量 $x$，有
\[
\lambda_3\le x^TAx\le \lambda_1.
\]

\textbf{【证明思路】} 已知 $3$ 阶实矩阵 $A$ 的特征值为 $\lambda_1\ge\lambda_2\ge\lambda_3$。先证明对角矩阵情形，再用实对称矩阵可正交对角化推广。

\textbf{【详细证明】} 当 $A=\diag(\lambda_1,\lambda_2,\lambda_3)$ 时，对任意三维单位向量 $x=(x_1,x_2,x_3)^T$，
\[
x^Tx=x_1^2+x_2^2+x_3^2=1,
\]
且
\[
x^TAx=\lambda_1x_1^2+\lambda_2x_2^2+\lambda_3x_3^2.
\]
由 $\lambda_1\ge\lambda_2\ge\lambda_3$ 得
\[
\lambda_3\le x^TAx\le \lambda_1.
\]

若 $A$ 为实对称矩阵，则存在正交矩阵 $Q$ 使
\[
Q^TAQ=\diag(\lambda_1,\lambda_2,\lambda_3).
\]
令 $x=Qy$，因 $Q$ 正交，$x^Tx=y^Ty=1$。于是
\[
x^TAx=y^TQ^TAQy
=\lambda_1y_1^2+\lambda_2y_2^2+\lambda_3y_3^2,
\]
由对角情形同理得
\[
\lambda_3\le x^TAx\le \lambda_1.
\]

\textbf{【结论】} 对角矩阵与实对称矩阵两种情形下结论均成立。
\end{answerbox}

\section{考前速背清单}

\begin{scorebox}[最后一遍]
\begin{enumerate}
\item 伴随矩阵先写 $AA^*=A^*A=|A|E$，不要直接交换矩阵乘法顺序。
\item 极大无关组看原矩阵主元列；线性表出要给具体系数。
\item 非齐次方程组通解必须写成“特解 + 齐次基础解系线性组合”。
\item 对角化题必须保证 $P$ 的列向量顺序与 $\Lambda$ 的对角元顺序一致。
\item 实对称矩阵正交对角化必须单位化特征向量，并写 $Q^TAQ=\Lambda$。
\item 二次型题要同时写标准形、规范形、正负惯性指数和秩。
\end{enumerate}
\end{scorebox}

\end{document}
